\documentclass[../../../include/open-logic-section]{subfiles}

\begin{document}

\olfileid{sth}{z}{sep}
\olsection{جداسازی}

با اصلی برای جایگزینی تصریح ساده‌انگارانه آغاز می‌کنیم:

\begin{axiom}[طرحوارهٔ جداسازی] برای هر فرمول~$\phi(x)$، عبارت زیر یک
اصل موضوع است: برای هر~$A$، مجموعهٔ~$\Setabs{x \in A}{\phi(x)}$ وجود دارد.
\end{axiom}

توجه کنید که این یک اصل موضوع منفرد نیست، بلکه \emph{طرحواره‌ای} از اصول
موضوع است. \emph{بی‌نهایت} اصل موضوع جداسازی وجود دارد؛ یکی برای هر
فرمول~$\phi(x)$. این طرحواره را به‌همین اندازه درست (و معمولاً) به صورت
زیر می‌نویسند:

\begin{defish}
برای هر فرمول~$\phi(x)$ که «$S$» در آن نیامده است، عبارت زیر یک اصل موضوع
است:
\[
	\forall A \exists S \forall x(x \in S \liff (\phi(x) \land x \in A)).
\]
\end{defish}

مطابق با قراردادی که در آغاز~\olref[sth][][]{part} یادآوری شد،
فرمول‌های~$\phi$ در اصول موضوع جداسازی می‌توانند پارامتر داشته
باشند.\footnote{برای توضیح معنای این سخن، به بحث بلافاصله پس از
\olref[sfr][infinite][induction]{natinductionschema} بنگرید.}

جداسازی بی‌درنگ با تصور انباشتی-تکراری‌ای که از مجموعه‌ها وصف کرده‌ایم
توجیه می‌شود. برای دیدن دلیل، فرض کنید~$A$ یک مجموعه باشد. پس~$A$ تا
مرحله‌ای مانند~$S$ تشکیل شده است (بنا بر~\stageshier). چون~$A$ در
مرحلهٔ~$S$ تشکیل شده است، همهٔ اعضای~$A$ پیش از مرحلهٔ~$S$ تشکیل شده‌اند
(بنا بر~\stagesacc). اکنون همهٔ مجموعه‌هایی را در نظر بگیرید که هم عضو
$A$ هستند و هم~$\phi$ را ارضا می‌کنند؛ آشکار است که همهٔ این مجموعه‌ها
نیز پیش از مرحلهٔ~$S$ تشکیل شده‌اند. بنابراین، آن‌ها نیز در مرحلهٔ~$S$
به صورت مجموعهٔ~$\Setabs{x \in A}{\phi(x)}$ تشکیل می‌شوند (بنا
بر~\stagesacc).

این اصل، برخلاف تصریح ساده‌انگارانه، از پارادوکس راسل اجتناب می‌کند. زیرا
نمی‌توانیم صرفاً وجود مجموعهٔ~$\Setabs{x}{x \notin x}$ را ادعا کنیم. در
عوض، با \emph{داشتن} مجموعه‌ای مانند~$A$، می‌توانیم وجود مجموعهٔ
$R_A = \Setabs{x \in A}{x \notin x}$ را ادعا کنیم. اما همهٔ آنچه از این
به دست می‌آید آن است که~$R_A \notin R_A$ و~$R_A \notin A$، و هیچ‌یک
چندان نگران‌کننده نیست.

بااین‌حال، جداسازی پیامدی بی‌درنگ و چشمگیر دارد:

\begin{thm}\ollabel{thm:NoUniversalSet}
هیچ مجموعهٔ \emph{جهانی‌ای} وجود ندارد؛ یعنی~$\Setabs{x}{x = x}$ وجود ندارد.
\end{thm}

\begin{proof}
برای برهان خلف، فرض کنید~$V$ مجموعه‌ای جهانی باشد. آنگاه بنا بر جداسازی،
$R = \Setabs{x \in V}{x \notin x} = \Setabs{x}{x \notin x}$ وجود دارد
و این با پارادوکس راسل در تناقض است.
\end{proof}

نبود مجموعه‌ای جهانی---و در واقع، بازبودن انتهای سلسله‌مراتب
مجموعه‌ها---یکی از بنیادی‌ترین اندیشه‌های پسِ تصور انباشتی-تکراری است.
پس ارزش دارد ببینیم که به‌طور شهودی می‌توانستیم از راهی دیگر نیز به آن
برسیم. یک مجموعهٔ جهانی باید خود !!a{element} خویش باشد. اما بنا بر تصور
انباشتی-تکراری ما، هر مجموعه برای نخستین بار در نخستین مرحلهٔ بلافاصله پس
از همهٔ !!{element}هایش در سلسله‌مراتب ظاهر می‌شود. ولی از این نتیجه
می‌شود که \emph{هیچ} مجموعه‌ای عضو خود نیست. زیرا هر مجموعهٔ خودعضو باید
نخستین بار بلافاصله پس از مرحله‌ای ظاهر شود که خودش برای نخستین بار در آن
ظاهر شده است، و این محال است. (در~\olref[sth][spine][rank]{defnsetrank}
خواهیم دید چگونه با استفاده از مفهوم «رتبهٔ» یک مجموعه این توضیح را
دقیق‌تر کنیم. البته برای این کار باید چند اصل موضوع دیگر نیز در اختیار
داشته باشیم.)

در ادامه چند پیامد دیگر جداسازی و گسترش‌مندی می‌آید.

\begin{prop}\ollabel{prop:emptyexists}
اگر مجموعه‌ای وجود داشته باشد، آنگاه~$\emptyset$ وجود دارد.
\end{prop}

\begin{proof}
اگر~$A$ مجموعه باشد، بنا بر جداسازی
$\emptyset = \Setabs{x \in A}{x \neq x}$ وجود دارد.
\end{proof}

\begin{prop}
برای هر دو مجموعهٔ~$A$ و~$B$، مجموعهٔ~$A \setminus B$ وجود دارد.
\end{prop}

\begin{proof}
بنا بر جداسازی، $A \setminus B = \Setabs{x \in A}{x \notin B}$ وجود دارد.
\end{proof}

همچنین معلوم می‌شود که اشتراک‌های (تقریباً) دلخواه وجود دارند:

\begin{prop}\ollabel{prop:intersectionsexist}
اگر~$A \neq \emptyset$، آنگاه
$\bigcap A = \Setabs{x}{(\forall y \in A)x \in y}$ وجود دارد.
\end{prop}

\begin{proof}
فرض کنید~$A \neq \emptyset$؛ پس~$c \in A$ برای یک~$c$. آنگاه
$\bigcap A = \Setabs{x}{(\forall y \in A)x \in y} =
\Setabs{x \in c}{(\forall y \in A)x \in y}$، که بنا بر جداسازی وجود دارد.
\end{proof}

البته به شرط~$A \neq \emptyset$ توجه کنید؛ زیرا~$\bigcap \emptyset$ به‌طور
تهی‌صدق مجموعهٔ جهانی خواهد بود، برخلاف~\olref{thm:NoUniversalSet}.

\end{document}
